%======================%
%     Nama Autoref     %
%======================%

% --- Teorema ---
\renewcommand{\theoremautorefname}{\Teorema}
\newcommand{\lemmaautorefname}{\Lema}
\newcommand{\definitionautorefname}{\Definisi}
\newcommand{\exampleautorefname}{\Contoh}

% --- Nama Level Bagian untuk \autoref ---
%
%     Penamaan untuk level Chapter hingga Subparagraph saat 
%     dirujuk
%     di dalam kalimat, contohnya lihat Bagian 2.5.
%
%     [NOTE] Rujukan level bagian hanya efektif untuk penomoran 
%            heading Multilevel (bernomor 1, 1.1, 1.1.1, 1.1.1.1)
%
\ifdef{\captionsindonesian}{\addto\captionsindonesian{
    \renewcommand{\chapterautorefname}{Bab}
    \ifdef{\chapter}{%
        \renewcommand{\sectionautorefname}{Subbab}
    }{%
        \renewcommand{\sectionautorefname}{Bagian}
    }%
    \renewcommand{\subsectionautorefname}{Bagian}
    \renewcommand{\subsubsectionautorefname}{Bagian}
    \renewcommand{\paragraphautorefname}{Paragraf}
    \renewcommand{\subparagraphautorefname}{Subparagraf}
}}{}

% --- Nama Sisipan ---
%
%     Penamaan untuk gambar, tabel, matematika, dan kode yang ada
%     di keterangan lampiran maupun dirujuk di dalam kalimat; 
%     contohnya lihat Gambar 2, lihat Tabel 3, Persamaan 4 
%     menunjukkan, lihat Kode 5.
%
\ifdef{\captionsindonesian}{\addto\captionsindonesian{
    \renewcommand{\figureautorefname}{Gambar}
    \renewcommand{\tableautorefname}{Tabel}
    \renewcommand{\equationautorefname}{Persamaan}
}}{}

\ifdef{\lstlistingname}{%
    \ifdef{\captionsindonesian}{\addto\captionsindonesian{
        \renewcommand{\lstlistingname}{Kode}
    }}{%
        \ifdef{\captionsamerican}{\addto\captionsamerican{
            \renewcommand{\lstlistingname}{Code}
        }}{%
            \ifdef{\captionsbritish}{\addto\captionsbritish{
                \renewcommand{\lstlistingname}{Code}
            }}{}%
        }%
    }%
}{}


